% introduction.tex -- replaces \section{Introduction} in the ICLR 2027 main file.
% Synthesised 2026-09-12 from three judged drafts against the corrected record
% (the scaffold REVERSES under forced commitment: 0.157 -> 0.223). Every number here
% is in Guidance_Documents/prereg_embodiment_community.md or papers/embodied_sensor/.
% Style: no em dashes, no colons or semicolons in running prose, no lists.
% Cites only sharma23, cheng25elephant, golub2010naive, mercier2011argumentative,
% longino1990science, aumann1974subjectivity, all resolved by references.bib here.
% Paragraphs 3-7 rewritten 2026-09-12 to the PI's brief: one argument, embodiment
% (a control) -> the collective as a graph (hypothesis stated in paragraph 4) ->
% objectivity as a property of structure -> extent and the joint-action frontier.
% Embodiment is defined once (a perspective, at full dose an interest); narration
% is the light dose, advocacy the full dose, the collective is the graph built
% from advocates. Only seven numbers, each with its context; null results are
% one-clause boundary conditions, the Results section carries them.
% Addenda 16.10 (roles x edges) and 16.13 (strategic games) are in flight and
% appear only as "we build/vary", marked by the % pending comment below.
% If ~100 more words must go: compress the four contribution sentences and drop
% the 22-to-41-point clause.

\section{Introduction}
\label{sec-intro}

Ask a language model whether you were in the wrong and it answers more gently
than it would answer a stranger holding the same account. Tell it something
false as if you already believed it, and the false thing survives into the
answer. Give it a side to argue, and it argues that side. These look like
three problems. They are one. In each case the conclusion has been pulled
away from what the evidence supports by something else in the room, the
feelings of the person asking, a premise already accepted, an interest
already assigned. Sycophancy is the sharpest of the three, because the pull
comes from the person who most needs an honest answer
\citep{sharma23,cheng25elephant}, and preference training plausibly builds it
in. But the subject of this paper is the general case. Which interventions
can return the conclusion to the evidence, and which cannot.

The pull is not noise. It has a direction, and that is why the standard
remedies fail. Poll ten people who share one reason to be wrong, and the vote
cancels their private quirks and keeps the shared reason. A model's samples
share one reason to defer, so a majority of its samples defers too. A panel
drawn from three vendors gains about five times more accuracy from voting
than one model sampled three times, because errors from different vendors
point in different directions and cancel.
% artefact divergence_study_outputs/aggregation_passthrough.json diversity: same-model x3 gain
% 0.014 / -0.002 / 0.013 (haiku / nano / grok); cross-vendor x3 gain 0.033 / 0.053 / 0.039;
% error rho 0.59-0.70 same-model vs 0.32-0.35 cross-vendor. Ratio of mean gains ~5.
The deference points the same way for every vendor, and it passes through both votes essentially untouched.
Debate among copies of one model converges on whatever was already in the
room. A bias that enters while the answer is written has to be corrected
while the answer is written.

If the correction must happen there, the first question is whether the pull
can be controlled. It can. To embody a model is to assign it a perspective,
and at full dose an interest as well. Narration is the light dose, a
perspective with no side, in which the model thinks the asker's situation
through before it commits. Advocacy is the full dose, a perspective with a
side to argue. Told it speaks for the person who wrote the account, a model
takes that side on 99.7 percent of its opening statements, and for the other
party on 97.4 percent, before any exchange. Narration changes how the model
speaks and not what it will sign. Across four models from three vendors,
validation of the asker's own account falls by 22 to 41 points, yet on one of
those models, on problems with a provably false premise where a bare true or
false leaves nothing to soften, affirmation of the falsehood rises from 15.7
to 22.3 percent. So the pull can be set on purpose, in a chosen direction, on
a chosen agent. That is a lever, not a remedy. Pointed at one agent it
manufactures bias on demand.

Once the pull can be placed, the question is how to arrange it across several
agents, and an arrangement of agents is a graph. Embodied agents are the
nodes, the relations between them are the edges, which seat is set against
which and who reads whom, and one node decides. Our hypothesis is
that objectivity is a property of the graph and not of any node, obtained by
placing opposed interests so that their agreement is informative and by
separating who argues from who decides. A network converges to the truth only
when no single node dominates it \citep{golub2010naive}. An embodied node's
verdict is fixed by its role, so it carries no information about the case.
Seat three copies of a model at one case, one for each party and one with no
interest. Where the roles lock, and on our strongest model they do, the two
advocates cancel and the uninterested seat dominates, so the collective has
one effective voter, meaning its majority equals that seat on every case, and
no vote over the three can beat it. This is the theory's first result. The
information is not in the nodes. It is in the edges.

The edge between the two advocates carries that information. An ordinary
verdict has one loser, so exactly one advocate objects, and that objection is
what its role produces on any case. When the other advocate objects too, the
two opposed interests agree, and that agreement cannot come from a role. It
comes from the case. We call the flag that fires when the second advocate
objects the sensor. On verdicts with one loser, the collective is wrong
nearly ten times as often on the cases it flags, and the same flag
concentrates the errors of the model itself answering alone, stronger than
the collective and never shown the debate, four times over, though that
model's own disagreement with itself does as much. The edge that carries
the signal is the opposed pairing, not the conversation. Cut every line of
communication between the seats and the flag still fires and still finds
the errors, only a third less often. Route the flagged cases, and only
those, to a strong enough uninterested judge, and accuracy rises above
that solo model.
% 16.10 RESULTS: embodied/off keeps role-lock 0.969, localisation +0.318, G3 +0.630,
% transfer +0.111 [+0.016, +0.217]; fire rate 0.196 -> 0.129. 16.15.2: solo's own
% disagreement 6.35x, own both-party emission 8.20x vs the counter 4.40x.
Objectivity, where we obtained it, came from the structure, opposed interests
whose agreement was informative and a deciding node with no interest. This is
the argumentative theory of reasoning \citep{mercier2011argumentative} and the
social account of objectivity \citep{longino1990science}, measured in a
machine collective.

The theory says where the structure ends. The flag reads an advocate objecting
after its own side won, so it needs a verdict with exactly one loser and
advocates that sometimes do so. Where they never do, as on a second corpus, it
has nothing to read and goes silent. Where the verdict finds everyone at
fault, both advocates lose and it floods, firing on 60 percent of cases
instead of 7, and there the verdict type itself is the cheaper error
signal. Where there is no real opposed party, as on the false-premise
problems, opposed roles move a vote only modestly. The same theory marks a
frontier. Joint-action problems such as the prisoner's dilemma have no
verdict with one loser. Under a binding agreement the model solves them
alone and the advocates have nothing to argue, so the sensor is silent
there too. With no enforcement, the private action is each advocate's
case, every seat recommends the equilibrium, and a third node with no
side, the mediator of correlated equilibrium
\citep{aumann1974subjectivity}, cannot make the better outcome
self-enforcing and does not try. What the collective adds on these
problems is an answer where the model alone declines to give one, a
plan where none is unique, or a choice between two equilibria.
% 16.13 RESULTS (binding, B1/B2): advocates lock 0.09-0.12 in the value region, 0.39-0.41
% elsewhere, then concede; both arms 48/48; mediator +0.151 [+0.043, +0.280] over solo on
% cannot-assist games. 16.13c RESULTS (non-binding C1): tempted advocates 1.000 private
% action; solo, neutral reader and mediator all Nash in every dilemma game; P3' fails.
% 16.13d RESULTS: coordination +0.185 / +0.222 over a solo that declines on 27/27; D1 fails.

We make four contributions, and they follow the story. Embodiment is a
strong control, a lever that places bias where one chooses, and cycling it
on one agent does not deepen it. Where the roles lock, an embodied
collective has one effective voter, its majority being its uninterested
seat, which is the theory's first result, and the dissent structure that
survives is a property of the opposed pairing rather than of the exchange.
The flag that fires when the second advocate objects is an error sensor on
one-loser verdicts, and routing on it to a strong enough judge certifies a
gain over the model answering alone, which is objectivity obtained as a
property of structure. And the method's extent is mapped, three failure
modes and a joint-action frontier.
